% Online supplement to:
%   A Comparability Protocol for Benchmarking Relational Database Access Layers in Express.js
%   in Express.js for PostgreSQL and MySQL
% Build: latexmk -pdf supplement.tex   (from paper/; tables in tables/)
\documentclass[11pt,a4paper]{article}
\usepackage[utf8]{inputenc}
\usepackage[T1]{fontenc}
\usepackage{booktabs}
\usepackage{array}
\usepackage{longtable}
\usepackage{adjustbox}
\usepackage{pgfplots}
\pgfplotsset{compat=1.18}
\usepackage[protrusion=true,expansion=false]{microtype}
\usepackage{url}
\usepackage{hyperref}
\hypersetup{colorlinks=true, linkcolor=blue, urlcolor=cyan,
  pdftitle={Online supplement: Relational Database Access-Layer Performance in Express.js}}

\newcommand{\tabledir}{tables/}
\renewcommand{\thetable}{S\arabic{table}}
\renewcommand{\thefigure}{S\arabic{figure}}

\title{Online supplement:\\
  A Comparability Protocol for Benchmarking Relational Database\\
  Access Layers in Express.js}
\author{Mateusz Miotk\\
  Faculty of Mathematics, Physics and Informatics,\\
  University of Gda\'nsk\\
  \texttt{mateusz.miotk@ug.edu.pl}}
\date{}

\begin{document}
\maketitle

\noindent This supplement carries measurement detail referenced from the main text. The complete revision artifact is archived as release v1.13.19, which the concept DOI \url{https://doi.org/10.5281/zenodo.21313858} resolves to; the Zenodo record also carries an immutable per-version DOI for that release. This release adds the specification-derived oracle, campaign-state preflight, a replacement corrected-state primary campaign, a second same-host RQ2 validation campaign, fresh secondary measurements, the external-audit dataset, and blind review packets. The revised primary tables are generated from the corrected-state campaign rather than inherited from the earlier baseline.

\smallskip
\noindent\emph{AI-assisted research tooling.} The harness, table generators, and statistical
estimators behind these tables were implemented with assistance from Claude (Anthropic), which also
proposed candidate analysis procedures that the author selected among, verified, and interpreted.
The models, affected components, the checks that independently validate the generated code, and the
limits of those checks are documented in the main text (Study Design, ``AI-assisted research
tooling''). Every table below reporting a measurement is regenerated from archived raw records by a
committed generator listed in \texttt{MANIFEST.md}; \texttt{MANIFEST.md} marks the authored
analytical tables, which report no measurement.

\section{Dispersion on MySQL}
\label{sup:cv-mysql}
Table~\ref{tab:cv_mysql} reports run-level throughput dispersion on MySQL. All five patterns use the one corrected-state campaign. Figure~\ref{fig:insert_dispersion} foregrounds all corrected insert replicates because a median and CV can hide mixture or tail structure. Historical insert samples from the defective reset are preserved in the archive but are not used for the revised ordering.

\input{\tabledir cv_mysql}

\input{\tabledir fig_insert_dispersion}

\section{Policy-selected documented round-trip counts}
\label{sup:query-counts}
Table~\ref{tab:query_counts} reports, for each layer, the number of SQL
statements the policy-selected documented deep fetch issues per request, captured from
server-side statement logs. The counts motivate the main text's observation
that round-trip count alone does not determine throughput: on the PostgreSQL
deep fetch single-query MikroORM is slowest, while two-query Drizzle outruns
four-query Prisma.

\input{\tabledir query_counts}

\section{Single-row insert throughput under default versus relaxed durability}
\label{sup:durability}
Table~\ref{tab:durability} gives the per-layer insert throughput under each
engine's default durability configuration (the paper's primary single-row-insert regime)
and under the symmetric relaxed regime (secondary, 5 replicates), with the
relaxed-to-default ratio quoted in the main text.

\input{\tabledir durability}

\section{Equal-CPU-budget sensitivity check}
\label{sup:equalcpu}
Table~\ref{tab:equalcpu} gives the deep-fetch throughput with the server
confined to 1, 2, or 4 cores (database and load generator pinned to disjoint
cores). The main text states the equal-CPU check as an operating condition; the nine values are
reported here.

\input{\tabledir equalcpu}

\section{Application-tier versus end-to-end CPU efficiency}
\label{sup:cpu}
Table~\ref{tab:cpu_efficiency} reports, per layer on the deep fetch,
the median application-process-tree CPU and database-server CPU, the combined
core-count, and throughput per application-CPU-second and per combined
(application$+$database) CPU-second. The application-only figure understates the
compute of layers that delegate more work to the database, but on the current
versions the application-only and combined figures largely agree, because no
layer draws more than about one application core: \texttt{mysql2} leads
application-tier efficiency by about $3.8\times$ over Prisma (2{,}301 vs.\
600~req per app-CPU-s) and by $3.4\times$ end-to-end (1{,}272 vs.\ 371~req per
combined-CPU-s). There is no CPU-for-throughput trade of the kind an earlier
version of this study reported.

\input{\tabledir cpu_efficiency}

\section{Coordinated-omission-corrected open-loop sweep}
\label{sup:openloop}
Table~\ref{tab:openloop} gives the full constant-arrival sweep (PostgreSQL,
five layers, offered rates 250--4{,}000 requests/s) behind the open-loop
paragraph of the main text: achieved throughput, corrected latency
percentiles measured from each request's \emph{intended} send time, and
timeout counts (10~s cap, clipped into the distribution).

\input{\tabledir openloop}

\section{Pool-size sensitivity}
\label{sup:poolsize}
The primary campaign fixes every layer's connection pool at 10 to isolate the
access layer from pool tuning. Table~\ref{tab:poolsize} varies the pool from 1 to
50 for one layer of each tier on the deep fetch (PostgreSQL, 50
connections). The deep-fetch ordering (\texttt{pg} $>$ Knex $>$ Prisma $>$ MikroORM) holds at every pool size of
two or more, and each layer's throughput saturates at a pool well below 50, so the fixed-pool choice
does not appear to drive the ranking. A pool of one is the exception and the one regime that
compresses the differences, because every layer then queues on connection acquisition: there
Prisma (870) overtakes Knex (741) on PostgreSQL.

\input{\tabledir poolsize}

\section{Transactional multi-statement write}
\label{sup:txn}
Beyond the single-row insert of the primary matrix, Table~\ref{tab:txn_write}
reports a transactional write, \texttt{POST /threads}, which inserts a post and
five comments in one transaction through each layer's transaction facility, for
one layer of each tier on PostgreSQL under default durability (median of five
runs, exploratory). The
full layer\,$\times$\,engine matrix for this pattern is left to future work. The
transaction commits once; its tighter clustering than the widest read spread
is consistent with amortizing per-commit work ($3.7\times$, \texttt{pg} 2{,}718 down to
Prisma 733), though the ordering does not simply track the single-row insert:
Prisma, mid-pack on the single-row PostgreSQL insert, falls to the bottom of the
transactional subset.

\input{\tabledir txn_write}

\section{Dispersion on PostgreSQL}
\label{sup:cv-pg}
Table~\ref{tab:cv} gives the full per-layer, per-pattern coefficient of variation
of throughput on PostgreSQL across the 25 repeated runs (maximum 6.4\%,
\texttt{pg}/insert), the companion to the MySQL table in Section~S1.

\input{\tabledir cv_all}

\section{Resource footprint}
\label{sup:resources}
Table~\ref{tab:resources} gives the per-layer application-process CPU (median \%
of one logical core) and peak resident memory on the deep fetch
(MySQL), the resource-accounting detail supporting the main text's operating conditions
(all layers $\approx100$--$110\%$ of one core; peak RSS 223--356~MB).

\input{\tabledir resources}

\section{Pinned versions}
\label{sup:versions}
Table~\ref{tab:versions} lists the pinned versions of every evaluated access layer
and tool, recorded from the lockfile on the measurement date.

\begin{table}[htbp]
  \centering
  \caption{Pinned versions of the evaluated access layers and tooling.}
  \label{tab:versions}
  \begin{tabular}{l l l l}
    \toprule
    Layer / tool & Version & Layer / tool & Version \\
    \midrule
    \texttt{pg}         & 8.22.0  & \texttt{sequelize}      & 6.37.8 \\
    \texttt{mysql2}     & 3.23.0  & \texttt{typeorm}        & 1.1.0  \\
    \texttt{knex}       & 3.3.0   & \texttt{objection}      & 3.1.5  \\
    \texttt{drizzle-orm}& 0.45.2  & \texttt{@mikro-orm/core}& 7.1.6  \\
    \texttt{@prisma/client} & 7.8.0 & \texttt{express}     & 5.2.1  \\
    \texttt{autocannon} & 8.0.0   & Node.js                 & 24.18.0 \\
    PostgreSQL          & 18.4    & MySQL                   & 9.7.1  \\
    \bottomrule
  \end{tabular}

  \vspace{2pt}
  \footnotesize Pinned to the latest stable release compatible with the harness
  adapter contract at the freeze date (2026-07-15); Sequelize's version-7 line is
  a prerelease, so the 6.x stable line is used. Prisma~7 is Rust-free and connects
  through an official JS driver adapter (\texttt{@prisma/adapter-pg} for PostgreSQL,
  \texttt{@prisma/adapter-mariadb} with the \texttt{mariadb} 3.5.3 driver for MySQL,
  as Prisma ships no \texttt{mysql2} adapter).
\end{table}

\section{Per-pattern detail: the read patterns}
\label{sup:reads}
The two cheapest read patterns show the smallest between-layer spread and are
summarized in the main text; their full throughput and p99 tables are given here.
Table~\ref{tab:point_read} reports the point read (primary-key lookup) and
Table~\ref{tab:range_scan} the keyset range scan, each by access layer and engine.
The two patterns foregrounded in the analysis (deep fetch and insert) are tabulated
in the main text; aggregation is reported in full below.

\input{\tabledir point_read}
\input{\tabledir range_scan}

\section{Common-SQL raw-path sensitivity contrast: full table}
\label{sup:sameplan}
Table~\ref{tab:sameplan} gives the full common-SQL raw-path sensitivity contrast on the deep fetch for both engines: throughput when every layer executes identical statements verified by server-side capture. The smaller residual spread is a compound sensitivity contrast, not a bound or causal decomposition.

\input{\tabledir sameplan}

\section{Per-pattern detail: aggregation}
\label{sup:aggregation}
Table~\ref{tab:aggregation} gives the full policy-selected aggregation-pattern throughput and p99 by access layer and engine. RQ3 finds aggregation the least layer-sensitive of the five primary patterns; the shared logical query avoids treating differing aggregate semantics as a performance result, but does not isolate mapping or protocol costs.

\input{\tabledir aggregation}

\section{Policy-selected documented deep-fetch choices per access layer}
\label{sup:adapter-choices}
Table~\ref{tab:adapter_choices} documents, for each access layer, the exact
eager-loading API used on the deep fetch, the number of SQL round trips it
issues, the documented alternative loading strategy we did \emph{not} use, and the
query-logging state. Every layer uses the path selected by the frozen-page policy;
query logging is disabled and verified on all eleven adapters, and no lifecycle
hooks or validation run on the read path. The six data-mapper and ORM layers split
both ways on the join-versus-select-in selected path: Sequelize, TypeORM, and MikroORM
use a single selected join, while Prisma and Objection use selected separate batched
queries. The loading-strategy sensitivity check reported in the main text
(Table~\ref{tab:altloading}) switches the join-selected layers to select-in where
that alternative is a byte-identical drop-in. Full per-adapter detail is in the
replication package's \texttt{METHODOLOGY.md}.

\input{\tabledir adapter_choices}

\section{Open-loop tail latency on MySQL}
\label{sup:openloop-mysql}
Table~\ref{tab:openloop_mysql} is the MySQL companion to the PostgreSQL open-loop
sweep (Supplement Table~S6): coordinated-omission-corrected p99 on the deep
fetch at a constant offered rate. The sub-saturation tail is flat and small on
both engines, and each layer collapses once the offered rate crosses its capacity,
so the high-load-tail ordering holds on MySQL as on PostgreSQL.

\input{\tabledir openloop_mysql}

\section{Alternative eager-loading sensitivity}
\label{sup:altloading}
Table~\ref{tab:altloading} reports the loading-strategy sensitivity check: for the
data-mapper ORMs whose alternative strategy is a byte-identical drop-in, policy-selected documented
deep-fetch throughput versus the alternative, on both engines. Both endpoints were
verified to return byte-identical responses before timing.

\input{\tabledir altloading}

\section{MySQL insert commit-path wait sensitivity}
\label{sup:waitevents}
Table~\ref{tab:waitevents} reports \texttt{performance\_schema} redo-log and
binary-log wait measurements around 4,000 valid inserts per displayed layer at
default durability. Waits are summed across ten concurrent requests, so aggregate
wait per insert may exceed request wall time and the ratio is not a causal time
fraction. Similar wait totals and the relaxed-durability control are consistent
with a shared commit-path contribution; they neither establish a performance
floor nor decompose engine, driver, and adapter costs.

\input{\tabledir waitevents}

\section{Post-restart environmental-robustness check}
\label{sup:postreboot}
Table~\ref{tab:postreboot} reports the environmental-robustness check: as the
campaign's final stage, after a full restart of both database engines (cold buffer
pools, fresh connections), a selected deep-fetch subset was re-measured against
the primary median. The three-layer order is unchanged on PostgreSQL (\texttt{pg}, Prisma, MikroORM)
and MySQL (\texttt{mysql2}, Prisma, MikroORM); absolute throughput is up to $16\%$ lower, a cold-buffer and run-to-run
drift the paper's relative analysis is designed to absorb. Because PostgreSQL and
MySQL cells are not measured simultaneously, this drift is non-uniform across engines
($0\%$ on \texttt{mysql2}, $16\%$ on \texttt{pg}, $16\%$ on MikroORM/PostgreSQL), so
the cross-backend ratios remain within-campaign estimates from sequential, randomized cells and can carry residual drift. Broad deep-fetch gaps exceed this check, but close aggregation and insert reversals are not treated as durable transfer evidence; the corrected-state whole-campaign sensitivity is reported separately.

\input{\tabledir postreboot}

\section{Utilization-controlled tail latency}
\label{sup:utilization}
Tables~\ref{tab:utilization} and~\ref{tab:utilization_mysql} give the
coordinated-omission-corrected p99 on the deep fetch when each layer is offered
50/70/85/95\% of its $\widehat C_{50}$ capacity proxy (the 25-run median
50-connection deep-fetch rate), replicated five times, on both engines. At the stable nominal 50\% target the tails contract to 2.3--4.9~ms on PostgreSQL and 1.9--5.5~ms on MySQL. Most 70\% cells remain small. This is consistent with capacity and queueing contributing materially to the wider equal-demand ladder; it neither identifies their share nor establishes a deployment-general sub-saturation bound. Re-expressing the loads against the
full-sweep maximum leaves the 50\% and 70\% bands below saturation (Table~\ref{tab:capacity_sensitivity}). This is the
capacity-normalized latency
companion to the high-load primary p99. The 50\% and most 70\% cells, which carry that
comparison, are stable across runs; the 85/95\% cells are increasingly noisy (p99
varies several-fold between runs as a layer approaches saturation), so those columns
are read as trend, not precise values.

\input{\tabledir utilization}
\input{\tabledir utilization_mysql}

\section{Longer-window tail sensitivity}
\label{sup:taillong}
The primary p99 is measured over a 12~s window, which yields only approximately
60 top-percentile observations in the slowest deep-fetch cells. Table~\ref{tab:taillong}
re-measures the same operating point over 60~s (ten runs per cell). The
12-versus-60-second p99 rank correlation is 1.00 on both PostgreSQL and MySQL;
the maximum cell-median shift is 5.6\% and 3.8\%, respectively, and maximum
run-level p99 CV is 7.3\% and 3.2\%. The PostgreSQL difference reflects native-tie
resolution as well as endpoint-specific tail-window sensitivity. These data support
the broad practical ordering in this campaign, not an exact-invariance claim or a
general rule that a 12-second window is sufficient.

\input{\tabledir taillong}

\section{Multi-worker scaling}
\label{sup:cluster}
Table~\ref{tab:cluster} scales a selected subset (three layers per engine) to 1, 2, and 4
Express workers (a shared-port \texttt{node} cluster), median of three runs. Because each process
uses about one core and gains nothing from more, aggregate throughput rises roughly in proportion
to workers over this range: on PostgreSQL the native driver leads at one worker
(\texttt{pg} 3{,}299 to Prisma's 1{,}117); its aggregate throughput then plateaus by two
workers (4{,}055, then 3{,}961 at four) --- whether from database contention, shared-host CPU
competition, or the fixed-connection generator is not instrumented here --- while Prisma, low per
process, rises 1{,}117 to 2{,}275 to 4{,}449 and reaches native-like aggregate
throughput at four workers. This is a bounded check, not a demonstration of general scalability:
it stops at four workers, measures no database-side saturation (DB CPU, connections, or wait
events), and on MySQL horizontal scaling \emph{widens} the gap --- the native \texttt{mysql2}
keeps scaling to 8{,}883~req/s at four workers while Prisma reaches only 2{,}427 (about
\texttt{mysql2} at one worker). Within these limits a layer's low per-process throughput may be
addressed by additional application processes until another bottleneck is reached, at a
proportional cost in cores --- a partial test of the single-host resource-competition concern.

\input{\tabledir cluster}

\section{Pool-size frontier on MySQL}
\label{sup:poolsize-mysql}
Table~\ref{tab:poolsize_mysql} is the MySQL companion to the PostgreSQL pool sweep
(Supplement Table~S7): deep-fetch throughput as the connection pool varies from 1 to
50, per layer. Each layer saturates at a pool well below 50 and the ordering is
preserved at every pool size, so the fixed pool of ten does not appear to drive the ranking on
either engine.

\input{\tabledir poolsize_mysql}

\section{Mixed read/write workload}
\label{sup:mixed}
Table~\ref{tab:mixed} interleaves the deep-fetch read with single-row inserts at
90:10 and 70:30 read:write, both engines, with a physical write reset per run. The
layer ordering matches the read ordering and is insensitive to the mix, because the
deep-fetch read dominates the cost; exploratory.

\input{\tabledir mixed}

\section{Concurrency sweep}
\label{sup:scaling}
Figure~\ref{fig:scaling} sweeps the deep fetch from 1 to 256 client
connections for one layer of each tier, both engines. The 50-connection
operating point used throughout sits above the knee for every layer on this
pattern; the ordering is load-robust above about 32 connections. This is the evidence behind the choice of
operating point.

\input{\tabledir fig_scaling}

\section{Per-layer cross-backend-stack ranks}
\label{sup:ranks}
Table~\ref{tab:ranks} reports PostgreSQL and MySQL ranks for all five patterns from the accepted full corrected-state campaign. The superseded mapping is preserved separately in \texttt{rq2-historical-provenance-comparison.json}; its ineligible MySQL cells are not treated as replication evidence. A separate clean same-host validation campaign remeasures the seven portable layers on the reviewer-prioritized deep-fetch, aggregation, and insert patterns. Between-campaign agreement is descriptive, not cross-host replication, and close reversals inside the predeclared \(\pm5\%\) remain exploratory.

\input{\tabledir ranks}

\section{Layer-by-backend-stack interaction magnitude}
\label{sup:interaction}
Table~\ref{tab:interaction} reports the corrected-state PostgreSQL$\div$MySQL throughput ratio per layer and pattern with a within-campaign 95\% paired-bootstrap interval. It quantifies heterogeneity between the two tested stacks; because the DBMS, official adapter, lower-level driver, and protocol can change together, it is not a pure engine effect.

\input{\tabledir interaction}

\section{Study factors}
\label{sup:factors}
Table~\ref{tab:factors} records which factors the benchmark holds constant, which
it varies, and which are uncontrolled on the single virtualized host; the
uncontrolled factors are why results are reported for this configuration rather
than as a general library ranking.
\input{\tabledir factors}

\section{Paired significance of the p99 ladder}
\label{sup:significance-p99}
Table~\ref{tab:significance_p99} is the p99 counterpart of the throughput
significance table (Supplement Table~S36): the same paired permutation, Wilcoxon,
geometric-mean-ratio, paired-bootstrap-CI, and dominance analysis is applied to each
adjacent layer's per-replicate deep-fetch p99 on both engines. Differences of about
1~ms are practically unresolved at the histogram resolution and are not promoted to
substantive rank claims. These resolved gaps quantify high-load p99 under fixed demand.
Together with the matched-utilization sensitivity they are consistent with a material
capacity-and-queueing contribution, but do not identify a mechanism or establish
different sub-saturation latency.

\input{\tabledir significance_p99}

\section{Extended methods (in the replication package)}
The reporting checklist, the benchmarking-pitfalls checklist, the detailed statistical methods, and
the measurement details are distributed with the replication package
(\texttt{notes/supplement-methods.pdf}) rather than reproduced here, to keep this supplement compact;
the numbered tables and figures above are unaffected.

\section{Artifact reproducibility summary}
\label{sup:reproducibility}
Table~\ref{tab:reproducibility} summarizes what is needed to reproduce the study and which
results are regenerated automatically; \texttt{REPRODUCE.md} in the replication package is
the single runnable entry point.

{\small
  \begin{longtable}{@{}p{0.24\textwidth} p{0.68\textwidth}@{}}
  \caption{Artifact reproducibility summary: version and identifier, requirements, how to
    run, time and resources, and which results are reproduced automatically.}
  \label{tab:reproducibility}\\
    \toprule
    Item & Detail \\
    \midrule
    Artifact \& version & Complete revision artifact v1.13.19: seed-specification oracle, corrected-state 90-cell primary campaign, 42-cell same-host validation campaign, fresh secondary measurements, external-audit dataset, and blind review packets. \\
    Code provenance & Git tag \texttt{v1.13.19} identifies the exact release commit, also recorded in the GitHub and Zenodo metadata. The package-lock SHA-256 and aggregate/per-file SHA-256 cover 42 measurement, admission, setup, state-control, schema, and adapter files. Exact 42-file source archives are preserved for both accepted campaigns and verify member-by-member against their embedded manifests. The accepted primary aggregate is \texttt{532d9c7ffe94\ldots d06cf228}; the independently ordered same-host validation aggregate is \texttt{d32a32b5b68a\ldots 134d3789a}. Full values are retained in the source manifests. \\
    Persistent identifier & Concept DOI \texttt{10.5281/zenodo.21313858}, resolving to release v1.13.19; the Zenodo record also carries an immutable per-version DOI. \\
    Software & Node.js 24.18.0, PostgreSQL 18.4, MySQL 9.7.1, \texttt{npm}. Reference path:
      conda user-space (\texttt{conda-forge}); the Docker path pins the same versions by digest but starts relaxed durability in a containerized topology, so it reproduces the workflow rather than the headline default-durability insert numbers. \\
    Hardware \& disk & Single Linux host (full host and version capture in
      Table~\ref{tab:versions}); $\approx$1~GB for the seeded databases. \\
    Deterministic seed & 2{,}000 authors / 100{,}000 posts / 1{,}000{,}000 comments from a
      fixed PRNG. \\
    Smoke test ($\approx$5--10~min) & \texttt{npm ci}; \texttt{npm run migrate \&\& npm run
      seed}; \texttt{npm run bench:quick}; \texttt{npm run verify:spec} (seed-derived expected results), \texttt{npm run verify:seed-parity} (cross-engine fixture values), \texttt{npm run verify:fixed} (differential equivalence), \texttt{npm run verify:campaign-state}, and \texttt{npm run verify:writes} (write-state
      correctness, must print ALL WRITES SEMANTICALLY CORRECT). \\
    Full campaign & Primary matrix: \texttt{npm run campaign:rq2} produces 90 cells with 25 runs each, followed by \texttt{npm run analyze:rq2}. The reduced same-host validation uses \texttt{npm run campaign:rq2-validation} and produces 42 cells with 25 runs each. Both commands fail closed on request errors or an incomplete roster and record pre-warm-up startup retries separately. Secondary experiments require additional hours and are invoked individually in \texttt{REPRODUCE.md}. \\
    Regenerate from raw & $\approx$minutes, \emph{no database}: the standalone no-database renderers rebuild their mapped tables and figures from \texttt{results/*.json}, then
      \texttt{npm run sync:tables} and \texttt{make}; \texttt{MANIFEST.md} explicitly marks
      authored tables, the one pre-generated statement-log table, and seven run-coupled renderers. \\
    Verification status & On 1 August 2026, an author-run isolated reconstruction from a clean export of the tagged commit verified 67/67 archived JSON hashes. Of the 58 committed tables, 48 carry a generator and 10 are authored analytical items with none; the documented no-database chain regenerated 42 byte-for-byte, while six generator-backed tables (S2, S6, S14, S19, S25 and Figure~S2) were byte-compared unchanged because their renderer is run-coupled or their input, server statement logs, is not archived, and the 10 authored items were byte-compared unchanged (\texttt{notes/reconstruction-2026-08-01.md}). The earlier 26 July run covered the 60 datasets listed at that date; seven were added to the manifest afterwards. The earlier author-run archive-isolated v1.12.9 check verified 35/35 inputs and 45/50 then-committed tables; its five presentation-only differences are named in \texttt{notes/clean-room-reproduction.md}. Neither check is independent reproduction or a re-execution of the current measurements. \\
    Entry point & \texttt{REPRODUCE.md} (smoke test, full matrix, archive-isolated author reconstruction from the tarball); the per-table data-and-generator map is in \texttt{MANIFEST.md}. \\
    \bottomrule
  \end{longtable}
}

{\footnotesize
  \begin{longtable}{@{}p{0.34\textwidth} p{0.13\textwidth} p{0.11\textwidth} c p{0.15\textwidth}@{}}
  \caption{Scope of every experiment: the access pattern(s), engine(s), number of
    access-layer implementations, and repeated runs each covers. The \emph{primary
    comparison} (throughput and closed-loop p99) is the full five-pattern, two-engine
    matrix; every other experiment is a limited condition, most restricted to the
    deep fetch, and each conclusion is scoped accordingly. ``Layers'' counts
    implementations (the nine access layers plus, where present, the two tuned native
    baselines); a selected subset is used where a caption says so. Engine
    entries name where the experiment was run (\texttt{pg} is PostgreSQL-only,
    \texttt{mysql2} MySQL-only).}
  \label{tab:scope}\\
    \toprule
    Experiment & Pattern(s) & Engines & Layers & Runs \\
    \midrule
    Primary matrix: throughput \& closed-loop p99 (main paper) & all five & PG, MySQL & 11 & 25 independent \\
    Same-SQL raw-path contrast (Table~\ref{tab:sameplan}) & deep fetch & PG, MySQL & 9 & median of 10 \\
    Capacity-denominator sensitivity (Table~\ref{tab:capacity_sensitivity}) & deep fetch & PG, MySQL & 8 & archived-data re-expression \\
    Open-loop CO-corrected tail (Tables~\ref{tab:openloop},~\ref{tab:openloop_mysql}) & deep fetch & PG, MySQL & 5 & rate sweep \\
    Utilization-controlled tail (Tables~\ref{tab:utilization},~\ref{tab:utilization_mysql}) & deep fetch & PG, MySQL & 8 & median of 5 \\
    Equal-CPU, 1/2/4 cores (Table~\ref{tab:equalcpu}), exploratory & deep fetch & PG & 3 & median of 3 \\
    Node-cluster multi-worker (Table~\ref{tab:cluster}) & deep fetch & PG, MySQL & 3 & median of 3 \\
    Pool-size sweep, 1--50 (Tables~\ref{tab:poolsize},~\ref{tab:poolsize_mysql}) & deep fetch & PG, MySQL & 4 & median of 3 \\
    Mixed read/write, 90:10 \& 70:30 (Table~\ref{tab:mixed}) & deep fetch, insert & PG, MySQL & 5 & median of 5 \\
    Relaxed vs default durability (Table~\ref{tab:durability}) & insert & PG, MySQL & 11 & 25 default plus 5 relaxed \\
    Wait-event flush decomposition (Table~\ref{tab:waitevents}) & insert & MySQL & 3 & 4000 inserts/layer \\
    Transactional multi-statement write (Table~\ref{tab:txn_write}) & \texttt{POST /threads} & PG & 5 & median of 5 \\
    Longer-window (60\,s) tail (Table~\ref{tab:taillong}) & deep fetch & PG, MySQL & 11 & 10 runs \\
    Canonical-constructor microbenchmark (Table~\ref{tab:canonicalization_cost}) & all read shapes & none & shared code & 30 blocks \\
    Post-restart cold cache (Table~\ref{tab:postreboot}) & deep fetch & PG, MySQL & 3 & vs primary median \\
    Alternative eager-loading (Table~\ref{tab:altloading}) & deep fetch & PG, MySQL & 3 & median of 5 \\
    Same-host validation campaign & deep fetch, aggregation, insert & PG, MySQL & 7 portable & 25 independent \\
    Specification-derived read oracle (unnumbered panel in the semantic-admission section) & all read methods & PG, MySQL & 9 per engine & 1,000 random/type + boundary/list + six fan-out fixtures \\
    External protocol audit (Table~\ref{tab:protocol_retro}) & 9 source benchmarks & n/a & 7 stages & 1 author-rater \\
    \bottomrule
  \end{longtable}
}

\section{Construct validity of the policy-selected documented treatment}
\label{sup:construct}
The \emph{policy-selected documented} rule is a mechanical, predeclared treatment-selection policy with no behavioral practitioner interpretation. It selects the first relations API on the version-compatible official page under the stated conflict rule. For auditability, the artifact preserves the exact response returned by the nearest available Wayback capture at or before the freeze for each official URL; the versioned documentation-snapshot manifest (\texttt{manifest.json}) records the committed HTML, capture timestamp, SHA-256, pinned version, selected API, alternatives, tie-break status, and decision basis; the same per-layer assignment appears in \texttt{METHODOLOGY.md} and Table~\ref{tab:adapter_choices}. A capture establishes page state only at its timestamp, not continuous lack of change through the freeze. This provides author-replay provenance, not evidence of production frequency, optimality, or independent agreement. The author applied every assignment; two result-blind selection reviews are prepared but pending. Table~\ref{tab:construct} records this per layer. Drizzle is the one borderline case: it exposes both a SQL-style query builder (the selected join) and a higher-level relational-query API that its documentation also presents, so we selected the builder join under that base-layer rule --- the page presents it as the core SQL-style layer --- and record the relational-query API as the documented alternative.
Where the documented alternative is a byte-identical drop-in, the loading-strategy sensitivity
check (Table~\ref{tab:altloading}) adds an empirical signal: for the join-selected ORMs (Sequelize,
MikroORM) the policy-selected documented path is the \emph{faster} of the two strategies, so a
layer's deep-fetch deficit does not appear to be an artifact of a poor selected strategy; for Objection the documented
default (batched select-in) is the slower option on MySQL, which we disclose rather than correct; the
Drizzle, Prisma, and TypeORM alternatives were not evaluated (TypeORM's errors on the ordered
deep fetch), itself a portability caveat. The construct is therefore a reproducible policy-defined treatment, not a practitioner persona and not a claim about performance-tuned expert usage or measured production frequency; RQ1 is scoped accordingly (main text, Threats to Validity).

\begin{table}[htbp]
  \centering
  \footnotesize
  \caption{Construct-validity record for the policy-selected documented deep-fetch treatment: the registered API, what the frozen page establishes or leaves ambiguous, and the empirical loading-strategy signal where the documented
    alternative is a byte-identical drop-in (Table~\ref{tab:altloading}).}
  \label{tab:construct}
  \begin{adjustbox}{max width=\textwidth}
  \begin{tabular}{@{}l l p{0.30\textwidth} p{0.30\textwidth}@{}}
    \toprule
    Layer & Selected API & Frozen-page evidence / ambiguity & Loading-strategy signal \\
    \midrule
    \texttt{pg}, \texttt{mysql2} (+tuned) & hand-written JOIN & No relation API; parameterized driver query path & n/a (hand-written SQL) \\
    \texttt{knex} & builder \texttt{.join()} & Builder join page; no relation abstraction & n/a (hand-written SQL) \\
    \texttt{drizzle} & builder join & Two co-prominent paths; documented-base-layer tie-break selected builder & alternative (relational-query API) not evaluated \\
    \texttt{prisma} & \texttt{include} & Relations page presents \texttt{include}; alternative recorded & alternative (join strategy) not evaluated \\
    \texttt{sequelize} & \texttt{include} (join) & Eager-loading page presents \texttt{include} & selected path \emph{faster} than select-in \\
    \texttt{typeorm} & \texttt{relations} & Find-options page presents \texttt{relations} & alternative errors on the ordered deep fetch (not evaluated) \\
    \texttt{objection} & \texttt{withGraphFetched} & Eager-methods page presents \texttt{withGraphFetched} & selected path (select-in) slower than join on MySQL (disclosed) \\
    \texttt{mikroorm} & \texttt{populate} (join) & Populate page presents \texttt{populate} & selected path \emph{faster} than select-in \\
    \bottomrule
  \end{tabular}
  \end{adjustbox}
\end{table}

\section{Results roadmap}
\label{sup:outcomes}
Table~\ref{tab:outcomes} classifies every result by analysis role: primary measurements,
secondary controls, or exploratory sensitivity evidence. It is the navigational companion to the
reproduction map in \texttt{experiments/MANIFEST.md}, which ties each table to its generator and
input data. The policy-selected documented configuration is the RQ1 estimand; common-SQL,
alternative-loading, and operating-point experiments are sensitivity analyses, not alternative
practitioner personas.

\input{\tabledir outcomes}

\section{Per-pattern capacity and utilization at the operating point}
\label{sup:scaling-patterns}
The operating-point protocol requires that capacity be identified so the fixed operating point can be
read against it (main text, Study Design). The deep-fetch concurrency sweep
(Section~\ref{sup:scaling}, Figure~\ref{fig:scaling}) does this in per-layer detail for one pattern;
review point 6.5 asks whether the other four patterns, reported only at the fixed 50-connection
demand, sit at comparable utilization. To answer it we swept every access layer over a connection
ladder (1, 4, 8, 16, 32, 50, 100, 200) on \emph{all five} patterns and both engines, medianed over
three runs per point after a warm-up, and for each layer computed its estimated saturating throughput (the ladder
maximum), the utilization at 50 connections (throughput at 50 connections divided by that maximum),
and the knee (the smallest connection count reaching 95\% of the maximum).

Table~\ref{tab:scaling_patterns} reports the result. The knee lies at or below 50 connections for
every pattern on both engines, so the fixed 50-connection operating point places all five patterns at
high utilization of their own capacity: cross-pattern comparisons of p99 and spread at that point are
therefore made at comparable (high) relative utilization, not at unknown or widely different
fractions of capacity. The binding constraint is the ten-connection pool shared by every pattern
(main text, Controls), which is why the knee clusters well below 50 connections regardless of the
pattern's absolute throughput. This extends the capacity-identification stage of the protocol from
the deep fetch to all five patterns; the equal-utilization open-loop tail
(Section~\ref{sup:utilization}) and the exploratory equal-compute check (Section~\ref{sup:equalcpu}),
which each need a per-layer saturating-throughput \emph{target}, remain demonstrated on the deep
fetch, the pattern on which the layer spread is largest.

\input{\tabledir scaling_patterns}

\section{Paired significance of the deep-fetch ranking}
\label{sup:significance-deep}
The main text foregrounds paired \emph{effect sizes} for the deep-fetch ranking --- geometric-mean
per-replicate throughput ratios and the fraction of replicates in which the faster layer wins ---
because they, not the significance tests, carry the ordering. Table~\ref{tab:significance} reports the
post-hoc adjacent-pair tests as a descriptive companion: since the ranking itself selects which pairs
are compared, the inference is conditional, so the permutation and Wilcoxon results are not treated as
confirmatory. The bootstrap intervals are within-campaign (main text, Study Design), not general over
deployment environments.

\input{\tabledir significance_deep_fetch}

\section{Application-tier CPU trade-off}
\label{sup:cpu-tradeoff}
Figure~\ref{fig:cpu_tradeoff} plots each layer's application-tier CPU against its throughput on the
deep fetch, the resource-accounting stage the main text lists among its operating conditions: every layer
draws about one application core, and the native driver leads on throughput and on requests per CPU-second (application-tier and
combined), consistent with its leaner SQL; these CPU shares, measured at unequal throughput, do not
identify where per-request work occurs.

\input{\tabledir fig_cpu_tradeoff}

\section{Retrospective application of the protocol to prior benchmarks}
\label{sup:protocol-retro}
Table~\ref{tab:protocol_retro} reports a structured audit of the nine benchmark sources retained by the scoping search. The codebook fixes five judgement codes and seven stage definitions. The machine-readable audit preserves all 63 judgements with source locators and paraphrased evidence, and its validator checks completeness. No benchmark was rerun, and \emph{not reported} is not treated as proof that a procedure never occurred. The author performed the coding, so the exercise demonstrates a bounded application of the protocol and the narrower conclusions it produces, not inter-rater reproducibility or independent validation. Result-blind forms for additional raters are shipped in the artifact reviewer packet.

\input{\tabledir protocol_retro}

\section{Specification-derived read admission and differential equivalence}
\label{sup:semantic-equivalence}
The mandatory read gate has two layers. First, \texttt{verify-spec.mjs} imports an executable seed specification that replays the fixed data PRNG and campaign fan-out fixture without querying either database or importing any adapter. It derives expected scalar values, field sets, ordering, graph membership, and author aggregates for 1,000 random post and author identifiers, explicit boundary and list cases, and all six fan-out fixtures. Database-generated timestamps are checked for presence and canonical ISO-8601 representation. The specification-derived panel below reports 73,080 passing expected-result checks. A separate cross-engine fixture gate compares every declared fan-out value and generated identifier, excluding only DBMS-generated timestamps; its 662-row representation is identical on PostgreSQL and MySQL (SHA-256 \texttt{402e0aebbbc9\ldots a99fe5ca}). The full hash is retained in the artifact. This supplies ground truth relative to the declared deterministic task for the tested inputs; it is finite evidence, not proof for arbitrary programs.

\input{\tabledir spec_oracle}

Second, \texttt{verify.mjs} checks 12 fixed probes and \texttt{verify-property.mjs} differentially compares the implementations over a separate fixed-seed random and edge-case sweep. Table~\ref{tab:semantic_equivalence} reports 60,800 adapter-versus-native comparisons with no divergence. Here the native result is an equivalence comparator, not the expected-result oracle. A common error could pass this layer but fail the seed-derived layer; conversely, mutual disagreement identifies a comparability failure even when one implementation matches a small expected sample.

\input{\tabledir semantic_equivalence}

\section{The five access patterns}
\label{sup:patterns}
Table~\ref{tab:patterns} lists the five HTTP endpoints of the benchmark workload, the access
pattern each realizes, and what each is intended to stress. The main text (Study Design) describes
the patterns in prose; this table is the compact reference.

\input{\tabledir patterns}

\section{Protocol compliance levels and their coverage}
\label{sup:protocol-compliance}
The main text (Study Design) defines the six protocol compliance levels and their mapping to the
seven protocol stages: M1 and M2 form the mandatory validity core, and R3--R7 supply five
claim-specific extensions. Table~\ref{tab:protocol_compliance}
records precisely which workload cells satisfy each in this experiment. The mandatory validity core
and capacity characterization hold across all five patterns on both engines; the three
extensions --- common-SQL sensitivity, utilization-controlled latency, and resource accounting --- are
demonstrated on the deep fetch, the widest-spread tested pattern, as a representative-point
instantiation rather than a claim of full-matrix coverage. The sixth, implementation-waste level is satisfied mechanically for the measured source by the static dead-work check; its stronger independent-human-review form is not: packets exist under \texttt{notes/reviewer-packets/}, but no external human review is claimed. Stating the levels and their coverage
explicitly is what makes the protocol reusable: a later study can report its own compliance level per
cell against the same ladder.

\input{\tabledir protocol_compliance}

\section{Predefined contrasts against the native-driver reference}
\label{sup:native-contrasts}
Selecting adjacent pairs \emph{after} seeing the ranking makes those tests conditional on the observed
ordering, so their $p$-values add little (main text, Measurement stability; the adjacent-pair table is
Table~\ref{tab:significance}). A predeclared contrast set avoids this. Table~\ref{tab:native_contrasts}
fixes the reference (the native driver) and the contrasts (each portable layer against it) in advance,
independently of the ranking, and reports for each layer on the deep fetch the geometric-mean paired
ratio of native-to-layer throughput, a within-campaign 95\% paired-bootstrap interval, and the paired
dominance. Every portable layer trails the native driver in all 25 replicates on both engines
(dominance $1.00$), and the ratios range from $1.49\times$ (\texttt{knex}, PostgreSQL) to $7.04\times$
(\texttt{mikroorm}, PostgreSQL) --- essentially the same spread (7.04 against the main text's 7.01, a different estimator), now as predeclared contrasts
rather than post-hoc pairwise tests. The contrasts are within-campaign, not deployment-general.

\input{\tabledir native_contrasts}

\section{Run-to-run p99 spread}
\label{sup:p99-spread}
Each cell's reported p99 is the median of 25 run-level p99 values, and each run-level p99 is itself
estimated from roughly the top 1\% of that run's requests, so a single run's p99 is noisy.
Figure~\ref{fig:p99_spread} plots all 25 per-run p99 values behind each deep-fetch cell (PostgreSQL)
with the median marked, so the run-to-run stability of the estimate is visible rather than summarized
by a single number: the spread is tight for the fast layers and wider for the slow ones, consistent
with the coarser p99 estimate at higher latency, and the medians preserve the reported ordering. The
load generator reports latency percentiles rather than raw per-request latencies, so a full
HDR-histogram reconstruction is left to future work; the median-of-run-level-p99 estimator, its
within-campaign bootstrap interval, and the longer-window (60~s) sensitivity check
(Section~\ref{sup:taillong}) are the mitigations used here.

\input{\tabledir fig_p99_spread}

\section{Components moved by the common-SQL raw-path sensitivity contrast}
\label{sup:samesql-components}
The common-SQL raw-path sensitivity contrast (main text, RQ1) changes more than SQL text. Replacing each layer's policy-selected relation-loading path with the same hand-written SQL through its raw facility changes several mechanisms at once. Table~\ref{tab:samesql_components} enumerates them. The
residual spread it leaves is a common-SQL raw-path contrast --- not a bound on the combined effect, and
not attributable to any single component, in particular not to SQL formulation or query count.

\input{\tabledir samesql_components}

\section{Deep-fetch tail: equal closed-loop demand and matched utilization}
\label{sup:tailregimes}
The high-load p99 under equal closed-loop demand (main text, Tail latency) and the p99 at matched
fractions of each layer's own capacity read the deep-fetch tail two ways. Table~\ref{tab:tail_regimes}
shows both per layer. At equal demand, p99 spans 18--108~ms on PostgreSQL
and 28--121~ms on MySQL; at the stable nominal 50\% target it contracts to
2.3--4.9~ms and 1.9--5.5~ms, respectively. Most 70\% cells also remain small.
This is consistent with a material capacity-and-queueing contribution but does not
identify a causal share or a deployment-general latency bound. The unstable
85--95\% cells are excluded from the contraction claim.

\input{\tabledir tail_regimes}

\section{Standalone cost of canonical response construction}
\label{sup:canonicalization}
The shared constructors enforce comparable output, but could themselves mask layer cost if they performed substantial graph work. Pattern by pattern, the point read copies seven declared fields and normalizes identifiers, views, the boolean, and timestamp; the range scan maps that same operation over the already ordered 20-row array without sorting; the deep-fetch graph path copies five post fields, one three-field author, and the already ordered comment array with each three-field author, while the flat-row variant performs the same field projection over rows already grouped by the adapter; and aggregation converts four scalar fields. The insert response does not use a graph constructor: each adapter returns only \texttt{\{id: Number\}}, and the out-of-band state check validates the persisted row. They only copy selected fields, normalize numbers, booleans, and ISO timestamps, and map arrays that the adapter has already grouped; they perform no query, join, regrouping, sorting, or cache access. Table~\ref{tab:canonicalization_cost} reports their isolated cost. Even the largest fixture is 0.146\% of the smallest corresponding primary-campaign HTTP p50; the ten-comment deep-fetch forms are 0.072\%. These are single-process development-host measurements and a scale comparison, not a decomposition of request time.

\input{\tabledir canonicalization_cost}

\section{Sensitivity of the matched-utilization capacity denominator}
\label{sup:capacity-sensitivity}
The utilization experiment used the 25-run median throughput at 50 connections as a
predefined capacity proxy. Table~\ref{tab:capacity_sensitivity} re-expresses the same offered
loads against a separate empirical denominator: the maximum observed in the full concurrency
sweep. The proxy is 93.2--102.1\% of that maximum across the sixteen layer--engine cells, so the
nominal 50\% and 70\% loads become 46.6--51.0\% and 65.2--71.5\%. This sensitivity analysis
changes no p99 observation and keeps both interpretation bands below saturation.

\input{\tabledir capacity_sensitivity}

\section{Same-host RQ2 campaign sensitivity}
\label{sup:rq2-cross-campaign}
Table~\ref{tab:rq2_cross_campaign} compares the accepted primary campaign with a second independently ordered same-host campaign for the seven portable layers on deep fetch, aggregation, and insert. Both campaigns contain 25 complete repetitions per cell and zero timed request failures. The table reports exact-rank agreement, rank correlation, and maximum median drift. A reversal enters the main RQ2 result only when it recurs in both campaigns, both stack-specific gaps exceed 5\%, and paired intervals exclude equality in opposite directions in both campaigns. This controls whole-campaign fragility on one host; it is not independent-host replication.

\input{\tabledir rq2_cross_campaign}

\section{Estimands and what the design does not identify}
\label{sup:estimands}
Table~\ref{tab:estimands} consolidates, for each research question, the quantity the design
estimates and the quantity it does \emph{not} identify. It is reference material for reading the
main text's claim limits in one place.

\input{\tabledir estimands}

\section{Protocol-development chronology}
\label{sup:protocol-chronology}
Table~\ref{tab:protocol_chronology} dates each protocol element against the measurement
campaigns, separating elements specified before any campaign from those added between campaigns
and those introduced after the accepted campaigns. The last group was applied retrospectively, to
archived state and the same deterministic seeded dataset, rather than gating the timed runs.
It supports the development-versus-validation distinction drawn in Section~4.

\input{\tabledir protocol_chronology}

\section{RQ2 leave-Prisma-out sensitivity}
\label{sup:rq2-leave-prisma-out}
Table~\ref{tab:rq2_leave_prisma_out} repeats the RQ2 reversal-promotion procedure on the six
portable layers other than Prisma, the one treatment whose lower-level MySQL driver differs. The
promotion rule is unchanged. No pair is promoted, while cross-stack rank agreement moves in both
directions: up on deep fetch and insert, down on aggregation. Holding out Prisma removes the driver substitution and the
treatment together and therefore identifies neither as the cause.

\input{\tabledir rq2_leave_prisma_out}

\section{Protocol normative detail}
\label{sup:protocol-normative}
This section carries the stage definitions, applicability limits, and compliance scheme
referenced from Section~4.

\paragraph{Admission versus diagnosis} Semantic admission asks two distinct questions. First, does an implementation satisfy expected results derived from the declared task and deterministic seed rather than from the timed native path? Second, are competing implementations mutually equivalent under the declared comparator? Failure of either read check, or a mutation reaching the wrong database state, excludes the cell. These are finite conformance tests, not proof of correctness for arbitrary inputs. Workload constraints disqualify only when preregistered; query count, joins, and loading strategy are otherwise measured properties of the treatment.

\paragraph{Applicability limits} the output comparator must match the output semantics: byte-identical comparison,
used here, is valid only when equivalent outputs are deterministic and canonically serializable, and
otherwise (unordered collections, floats, timestamps, nondeterministic identifiers) needs a semantic
comparator (set equality, tolerance, canonicalization). The experiment's outputs are deterministic and
canonically constructed, so byte-equality is valid here.

\paragraph{Compliance levels} M1 semantic admission and M2 treatment definition form the mandatory core because, without a declared task result and a stable configured treatment, there is no common comparison estimand. The remaining stages license additional interpretations without deciding whether a fixed-load result exists. Capacity characterization locates that result against a throughput knee. The common-SQL extension licenses a compound sensitivity contrast; utilization control licenses comparison at stated fractions of estimated capacity; resource accounting licenses a compute-and-memory reading or a labelled equal-budget sensitivity. Implementation-waste evidence records evidence against one class of avoidable work on the timed
path, not its absence. Here M1--M2 cover the full matrix, R4 every pattern, R3/R5/R6 the deep fetch; R7 is met for the measured source by the static dead-work check, its runtime companion covering later campaigns only. Omitting a recommended stage narrows the supported claim without invalidating the reported numbers: Partial coverage of a recommended stage narrows what it licenses without withdrawing the admission evidence that licenses the comparisons.

\smallskip
Supplement Table~S50 records how each control changed this experiment; that establishes consequence here, not universal necessity. Supplement Table~S47 consolidates, per research question, the quantity estimated and the quantity not identified. Supplement Table~S37 separately applies the codebook to nine external benchmark sources. The machine-readable audit preserves 63 source-located judgements, but one author performed the coding, so it demonstrates bounded protocol application, not inter-rater reproducibility. The record is structured for a second rater: additional codings attach alongside the author's, and a committed scorer (\texttt{npm run audit:agreement}) reports pooled percent agreement, Cohen's kappa with a bootstrap interval, and every disagreeing cell. Per-stage slices are $n{=}9$ with chance agreement up to $0.80$, so per-stage kappa is flagged unreliable and the pooled figure is the headline. Result-blind forms for three optional independent exercises ship under \texttt{notes/reviewer-packets}; all remain pending, and the scorer reports agreement as not computable while one rater is recorded, so an unperformed exercise cannot be counted as validation.




\section{Protocol-to-interpretation mapping}
\label{sup:protocol-mapping}
Table~\ref{tab:protocol_mapping} records, per protocol stage, the concrete consequence it had
for this study's reported interpretation. It is authored analytical reference material rather
than a measurement table.

\input{\tabledir protocol_mapping}



\section{Prior-work coverage}
\label{sup:prior-art}
Table~\ref{tab:prior_art} positions prior work along the four coverage axes this experiment
spans: taxonomy breadth, engine coverage, joint throughput/tail reporting, and vendor
independence. Each source covers at most two or three.

\begin{table}[htbp]
  \centering
  \small
  \caption{Prior work on relational-database access-layer performance,
    positioned along four \emph{coverage} axes this experiment spans; the paper's contribution
    is the comparability protocol of the main text, not this coverage. The main text cites each source by number.
    ``Products covered'' aggregates native drivers, query builders, and ORMs
    into one count; ``none'' marks studies that vary the engine or
    framework but not the access layer. \emph{Vendor involvement} marks maintainer-authored
    studies --- a conflict-of-interest property of the source, not a methodological guarantee that
    a non-vendor-authored benchmark is unbiased.}
  \label{tab:prior_art}
  \begin{adjustbox}{max width=\textwidth}
  \begin{tabular}{@{}llllc@{}}
    \toprule
    Study & Products covered & Engines & Metrics & Vendor involvement \\
    \midrule
    Colley et al.\ (2018) & one ORM vs. hand-written SQL example & SQL Server & query latency & None \\
    Procedia CS (2025)$^\dagger$ & Prisma vs.\ raw SQL & PostgreSQL & latency, CPU & None \\
    JCCT (2025) & 3 ORMs & PostgreSQL & latency, memory, CPU & None \\
    JCSI (2025) & 3 ORMs & PostgreSQL & per-stage latency & None \\
    Node-Bench (2023) & 5 access products (+11 frameworks) & MySQL & throughput, latency & None \\
    Salunke and Ouda (2024) & none (engine only) & PostgreSQL + MySQL & throughput, latency & None \\
    Drizzle suite & 9 products & PostgreSQL & throughput, latency, P95 & Vendor \\
    Prisma suite & 3 ORMs & PostgreSQL & median latency only & Vendor \\
    IMDBench & 4 ORMs + driver & PostgreSQL & throughput, latency & Vendor \\
    \midrule
    \textbf{This study} & \textbf{9 policy-selected documented (+2 tuned)} & \textbf{PostgreSQL + MySQL} & \textbf{throughput + response-time p99} & \textbf{None} \\
    \bottomrule
  \end{tabular}
  \end{adjustbox}
  \vspace{4pt}
  {\raggedright\footnotesize $^\dagger$Cited for its methodology only; its absolute speedup figures are not comparable to ours given the differing versions, hardware, workload, and harness.\par}
\end{table}

\section{Insert throughput and tail}
\label{sup:insert}
Table~\ref{tab:write} reports single-row insert throughput and response-time p99 per layer and
engine under each engine's default durability, with within-campaign bootstrap intervals. The
MySQL insert cells are the noisiest in the study, which is why no fine insert ordering is claimed.

\input{\tabledir write}

\section{Result-graph breadth sweep}
\label{sup:fanout}
Table~\ref{tab:fanout} varies the deep fetch's result-graph breadth from 0 to 500 child rows and
reports each layer's throughput multiplier relative to its engine's native driver. The multipliers
are near-flat on MySQL but not on PostgreSQL, where the native-relative spread widens from
$6.0\times$ to $11.7\times$: MikroORM and Sequelize fall further behind the native driver as
breadth grows, while Knex and Drizzle close on it. Breadth is the only dimension varied, so the
sweep bounds neither graph depth, nor schema, nor query mix, and three runs per cell make it
exploratory rather than an estimate.

\input{\tabledir fanout}

\section{Multiplicity and exchangeability in the paired analyses}
\label{sup:multiplicity}

\paragraph{Multiplicity across the adjacent-pair families}
Two further test families are reported besides the RQ2 promotion screen and its correction
(Section~\ref{sup:rq2-multiplicity}): the seven adjacent throughput
pairs on MySQL (Table~\ref{tab:significance}) and the fourteen adjacent p99 pairs, seven per
engine (Table~\ref{tab:significance_p99}). Both follow the \emph{observed} ranking rather than a
preregistered one, so they are secondary descriptive robustness checks and not a confirmatory
family. As a conservatism check we apply a Bonferroni correction: at $\alpha=0.05$ the
per-comparison threshold is $0.05/7=7.1\times10^{-3}$ for the throughput family,
$0.05/14=3.6\times10^{-3}$ for the p99 family, and $0.05/21=2.4\times10^{-3}$ if the two are
pooled. All seven throughput comparisons and thirteen of the fourteen p99 comparisons sit at the
permutation resolution floor $1/(B{+}1)=5.0\times10^{-5}$ with $B=20{,}000$ and survive every
threshold. One does not: the MySQL Sequelize-against-Objection p99 step has $p=0.0223$, which
clears an uncorrected $\alpha=0.05$ but fails all three corrected thresholds. That adjacent step
is therefore not resolved once multiplicity is accounted for, and we do not read an ordering from
it. We state this rather than omit it because the
absence of a multiplicity policy, not its outcome, is what a reader cannot check. The floor also
means the tests bound the $p$-value from above and do not resolve differences among the pairs:
a comparison at the floor is not evidence that one adjacent step is more separated than another,
for which the ratio and its interval are the informative quantities.

\paragraph{Exchangeability of the replicate blocks}
The paired analyses treat the 25 replicates as exchangeable blocks: within each replicate every
layer is measured once, from the same seeded identifier distribution, and the layer order is
randomized. The permutation tests permute layer labels within blocks, which is valid under the null
if block position carries no systematic effect. That assumption is checked directly only for the
write path, where an explicit forward/reverse order-invariance run is reported; for the read
workloads it is supported by design (randomized order, per-replicate reset, handler drain between
cells) and by the low within-campaign dispersion (deep-fetch CV $\leq4.6\%$), but no separate
block-order or time-trend diagnostic is displayed. Database and host state persist across a
campaign, so a slow drift shared by all layers would cancel in the paired contrasts while a
layer-specific drift would not. The paired intervals and tests therefore rest on an unchecked
exchangeability assumption for the read workloads, which we state as a limitation rather than
treat as established.

\section{Effect sizes for the promoted cross-stack reversals}
\label{sup:rq2-effects}
Table~\ref{tab:rq2_promoted_effects} sizes the four promoted reversals that the main text's RQ2
answer rests on. Each is a paired within-campaign ratio with a 95\% bootstrap interval on each
stack, so the reversal is visible as the two intervals falling on opposite sides of parity rather
than as a rank statement alone. The largest is the insert reversal (PostgreSQL $1.423$ against
MySQL $0.723$) and the smallest the deep-fetch pair against Objection ($1.155$ against $0.754$).
The MySQL insert interval is the widest in the table, consistent with that cell's $13.6\%$
within-campaign dispersion, which is why no fine MySQL insert ordering is claimed anywhere.
These intervals are within-campaign: they cover replicate-to-replicate variation on one host and
one software freeze, not cross-host or cross-version transfer.

\input{\tabledir rq2_promoted_effects}

\section{Familywise correction across the RQ2 screen}
\label{sup:rq2-multiplicity}
The promotion rule is applied to 21 layer pairs on each of three patterns, 63 comparisons per
campaign, without a familywise correction, and the main text states that recurrence across two
same-host campaigns does not bound familywise error. Table~\ref{tab:rq2_multiplicity} supplies the
correction that disclosure leaves open. The corrected hypotheses are the necessary
components of what the rule claims, which is a conjunction: the pair's ratio differs from~1 on both stacks, in opposite
directions, by more than the margin. Testing the cross-stack interaction instead would certify a
hypothesis the rule already implies, since two intervals on opposite sides of parity entail a
non-zero interaction, and so would not cover the selection. Each stack's per-replicate log ratios are
therefore tested against zero by two-sided sign-flip permutation, the two are combined by
intersection-union (the larger of the two $p$-values, the correct combination for a conjunction;
using two-sided components for a directional claim is conservative), and
the result is Holm-corrected across all 63 pair-pattern screens in a campaign. Holm retains 53 of 63 in the primary campaign
and 50 of 63 in the validation campaign. Applying surrogate direction and margin filters to the
survivors, computed on geometric means rather than the rule's medians, leaves exactly the four
pairs the screen promoted, and in the validation campaign those four plus one that the recurrence
requirement excludes; that agreement is descriptive concordance between two estimators, not
evidence that the promotion decision itself was corrected. The two estimators are close here: for
all four promoted pairs on both stacks and in both campaigns they agree within $2.2\%$, and each
pair's median-ratio bootstrap interval excludes the $\pm5\%$ band on both stacks in both
campaigns. The sign-flip null here is
symmetry of the per-replicate log ratios about zero, a different and equally unchecked assumption
from the within-block label exchangeability of Section~\ref{sup:multiplicity}. What this does and
does not settle is worth stating plainly: it shows that the per-stack nonzero-effect evidence for the four promoted pairs survives Holm
correction,
and it leaves untouched the two limitations the main text already reports: the $5\%$ margin was
fixed before both campaigns, but the composite direction-and-margin rule that uses it was specified
after they completed, and the two campaigns share a host and the same seeded draws. Correction cannot make a post-hoc screen prospective.

\input{\tabledir rq2_multiplicity}

\section{Equivalence assessment for the closest adjacent pair}
\label{sup:equivalence}
Table~\ref{tab:tost_closest} reports the $\pm5\%$ equivalence assessment the main text refers to.
Three properties bound how it may be read. First, the assessment is the confidence-interval form of
TOST: it asks whether the $90\%$ paired bootstrap interval of the geometric-mean ratio lies inside
the equivalence bounds, and reports no separate one-sided test statistics. Second, a negative result
means equivalence within the margin was \emph{not established}; it does not reject equivalence, and
here neither interval lies wholly outside the margin either, since both straddle the upper bound.
Third, the pair is selected from the observed ordering as the smallest adjacent ratio among the eight
policy-selected layers on that engine, so the selection is post hoc and the assessment is a sensitivity analysis
rather than confirmatory evidence. The margin shares that status: $\pm5\%$ is author-selected and
application-dependent rather than a universal threshold, and although it was fixed on 2026-07-13,
before the accepted campaigns, pilot data already existed at that date, so it is not a preregistered
bound. The tuned baselines are excluded because the paper defines them as disclosed reference points
rather than treatments; the closest adjacent pair including them would be
\texttt{pg-tuned} versus \texttt{pg}, which does satisfy the margin.

\input{\tabledir tost_closest}

\end{document}
